
\documentclass{article}
\usepackage{colortbl}
\usepackage{makecell}
\usepackage{multirow}
\usepackage{supertabular}

\begin{document}

\newcounter{utterance}

\centering \large Interaction Transcript for game `ifeval', experiment `ifeval', episode 526 with reflexion{-}Qwen3.5{-}4B.
\vspace{24pt}

{ \footnotesize  \setcounter{utterance}{1}
\setlength{\tabcolsep}{0pt}
\begin{supertabular}{c@{$\;$}|p{.15\linewidth}@{}p{.15\linewidth}p{.15\linewidth}p{.15\linewidth}p{.15\linewidth}p{.15\linewidth}}
    \# & \multicolumn{2}{c}{Player} && \multicolumn{2}{c}{Game Master} \\
    \hline

    \theutterance \stepcounter{utterance}  
    & & & \multicolumn{4}{p{0.6\linewidth}}{
        \cellcolor[rgb]{0.9,0.9,0.9}{
            \makecell[{{p{\linewidth}}}]{
                \texttt{\tiny{[P1$\langle$GM]}}
                \texttt{What are the best things to do in Rochester, New York? Can you write them as a song and use music notation in your response? Make sure to include the keywords 'festival' and 'river'.} \\
            }
        }
    }
     \\ \\

    \theutterance \stepcounter{utterance}  
    & \multicolumn{4}{p{0.6\linewidth}}{
        \cellcolor[rgb]{0.561,0.737,0.561}{
            \makecell[{{p{\linewidth}}}]{
                \texttt{\tiny{[P1$\rangle$GM]}}
                \texttt{Here is a song about the best things to do in Rochester, New York, featuring the keywords **festival** and **river**, accompanied by music notation.} \\
\\ 
\texttt{**Title:** The City of Roses \& The Genesee} \\
\texttt{**Tempo:** Moderate, upbeat folk{-}rock} \\
\texttt{**Style:** Acoustic guitar with a driving bassline} \\
\\ 
\texttt{\#\#\# Lyrics \& Notation} \\
\\ 
\texttt{**(Intro)**} \\
\texttt{| C | Am | F | G |} \\
\texttt{| :{-}{-}{-} | :{-}{-}{-} | :{-}{-}{-} | :{-}{-}{-} |} \\
\texttt{| *C* | *Am* | *F* | *G* |} \\
\texttt{| *C* | *Am* | *F* | *G* |} \\
\\ 
\texttt{**(Verse 1)**} \\
\texttt{| C | Am | F | G |} \\
\texttt{| :{-}{-}{-} | :{-}{-}{-} | :{-}{-}{-} | :{-}{-}{-} |} \\
\texttt{| The sun comes up on the **River** wide, |} \\
\texttt{| Where the water flows by the Genesee side. |} \\
\texttt{| We walk the streets where the history lives, |} \\
\texttt{| And the old brick buildings bring us to our lives. |} \\
\\ 
\texttt{**(Chorus)**} \\
\texttt{| F | G | C | Am |} \\
\texttt{| :{-}{-}{-} | :{-}{-}{-} | :{-}{-}{-} | :{-}{-}{-} |} \\
\texttt{| Oh, Rochester, you are the place to be, |} \\
\texttt{| Where the music plays and the people see. |} \\
\texttt{| Join the crowd at the **festival** bright, |} \\
\texttt{| Under the stars of the northern light. |} \\
\texttt{| (F) (G) (C) (Am)} \\
\\ 
\texttt{**(Verse 2)**} \\
\texttt{| C | Am | F | G |} \\
\texttt{| :{-}{-}{-} | :{-}{-}{-} | :{-}{-}{-} | :{-}{-}{-} |} \\
\texttt{| Go to the museum where the art is deep, |} \\
\texttt{| Or take a ride on the train to sleep. |} \\
\texttt{| The University grounds are green and tall, |} \\
\texttt{| A story of learning that we all recall. |} \\
\\ 
\texttt{**(Chorus)**} \\
\texttt{| F | G | C | Am |} \\
\texttt{| :{-}{-}{-} | :{-}{-}{-} | :{-}{-}{-} | :{-}{-}{-} |} \\
\texttt{| Oh, Rochester, you are the place to be, |} \\
\texttt{| Where the music plays and the people see. |} \\
\texttt{| Join the crowd at the **festival** bright, |} \\
\texttt{| Under the stars of the northern light. |} \\
\texttt{| (F) (G) (C) (Am)} \\
\\ 
\texttt{**(Bridge)**} \\
\texttt{| Em | C | G | C |} \\
\texttt{| :{-}{-}{-} | :{-}{-}{-} | :{-}{-}{-} | :{-}{-}{-} |} \\
\texttt{| From the hills of the west to the east, |} \\
\texttt{| We find our joy in this city's feast. |} \\
\texttt{| (Em) (C) (G) (C)} \\
\\ 
\texttt{**(Outro)**} \\
\texttt{| F | G | C | Am |} \\
\texttt{| :{-}{-}{-} | :{-}{-}{-} | :{-}{-}{-} | :{-}{-}{-} |} \\
\texttt{| By the **River** we sing, by the **festival** we go, |} \\
\texttt{| Rochester, Rochester, oh, what a place to be! |} \\
\texttt{| (F) (G) (C) (Am)} \\
\texttt{| (F) (G) (C) (Am)} \\
\\ 
\texttt{***} \\
\\ 
\texttt{\#\#\# Musical Notation (Tablature \& Chord Guide)} \\
\\ 
\texttt{Since I cannot generate an image file directly, here is the structured musical notation using standard chord symbols and a rhythmic guide to help you play or visualize the song.} \\
\\ 
\texttt{**Key:** C Major} \\
\texttt{**Time Signature:** 4/4} \\
\\ 
\texttt{```text} \\
\texttt{{[}Verse 1{]}} \\
\texttt{C  Am  F  G} \\
\texttt{| C  Am  F  G |} \\
\texttt{C  Am  F  G} \\
\texttt{| C  Am  F  G |} \\
\\ 
\texttt{{[}Chorus{]}} \\
\texttt{F  G  C  Am} \\
\texttt{| F  G  C  Am |} \\
\texttt{F  G  C  Am} \\
\texttt{| F  G  C  Am |} \\
\\ 
\texttt{{[}Bridge{]}} \\
\texttt{Em  C  G  C} \\
\texttt{| Em  C  G  C |} \\
\texttt{Em  C  G  C} \\
\texttt{| Em  C  G  C |} \\
\\ 
\texttt{{[}Outro{]}} \\
\texttt{F  G  C  Am} \\
\texttt{| F  G  C  Am |} \\
\texttt{F  G  C  Am} \\
\texttt{| F  G  C  Am |} \\
\texttt{```} \\
\\ 
\texttt{**Performance Notes:**} \\
\texttt{*   **Rhythm:** Play the chords with a steady "strum{-}strum{-}strum{-}strum" pattern.} \\
\texttt{*   **Dynamics:** Start the verse softly, build energy in the chorus, and keep the bridge gentle before the final energetic outro.} \\
\texttt{*   **Keywords:** Ensure the words "**River**" and "**festival**" are sung with emphasis on the downbeat of their respective lines.} \\
            }
        }
    }
    & & \\ \\

    \theutterance \stepcounter{utterance}  
    & & & \multicolumn{2}{p{0.3\linewidth}}{
        \cellcolor[rgb]{0.9,0.9,0.9}{
            \makecell[{{p{\linewidth}}}]{
                \texttt{\tiny{[GM$|$GM]}}
                {'keywords:existence': {'keywords': ['festival', 'river']}}
            }
        }
    }
    & & \\ \\

    \theutterance \stepcounter{utterance}  
    & & & \multicolumn{2}{p{0.3\linewidth}}{
        \cellcolor[rgb]{0.9,0.9,0.9}{
            \makecell[{{p{\linewidth}}}]{
                \texttt{\tiny{[GM$|$GM]}}
                \texttt{game\_result = WIN} \\
            }
        }
    }
    & & \\ \\

\end{supertabular}
}

\end{document}
